% ===================================================================
%  FIN DU LIVRET — feuillets 93 a 96 : TABLE DES MATIERES, acheve
%  d'imprimer, 4e de couverture.
%
%  Feuillet 93 : table des matieres, non foliotee (verifie sur le
%  scan : ni haut ni bas de page). Feuillet 94 : page blanche (verso
%  vierge). Feuillet 95 : acheve d'imprimer (« 44872. --- Bordeaux.
%  --- Imp. G. Delmas, rue Saint-Christoly, 10. »), seul sur une page
%  autrement blanche. Feuillet 96 : 4e de couverture, publicite pour
%  la « Collection Delmas », en FRANCAIS (non traduite) — le cadre
%  ornemental du fac-simile n'est pas reproduit, seul le texte l'est.
%
%  Le facsimile francais est une prise de vue, contraste inegal ; le
%  feuillet 96 a ete regenere a 2600 px (outils/.lekto) pour rester
%  lisible. Un mot y reste incertain malgre plusieurs agrandissements
%  successifs : le patronyme du feuillet 96, transcrit « Veillet-
%  Aboison » d'apres la graphie la plus probable.
%
%  Corps et interlettrages du feuillet 96 : PROVISOIRES, choisis pour
%  que chaque ligne relevee tienne dans \VUtexteLargeur sans Overfull.
%  A reprendre par outils/apparat.py comme le reste de l'apparat.
% ===================================================================

% --- Feuillet 93 (table des matieres, non foliote) -------------------
\begin{VUpage}[93]{}
%%K t99-tit-1 sub
\VUsaut{3.0mm}
\VUtitre{15.0pt}{60}{TABLE DES MATI\`ERES}
\VUsaut{2.4mm}
\VUfilet{20mm}
\VUsaut{8.0mm}

{\fontsize{10.2pt}{12.2pt}\selectfont
%%K t99-01-1 p
\noindent \textsc{Avertissement} \dotfill 3

%%K t99-02-1 p
\noindent \VUgras{Tableau 1.} --- L'École. --- Le Lycée. --- La Classe \dotfill 7

%%K t99-03-1 p
\noindent \VUgras{Tableau 2.} --- Le Corps humain. --- La Récréation. ---\nl
Les Jeux \dotfill 12

%%K t99-04-1 p
\noindent \VUgras{Tableau 3.} --- L'Enfance (Un Baptême au village) \dotfill 17

%%K t99-05-1 p
\noindent La Jeunesse (Une Fête publique) \dotfill 20

%%K t99-06-1 p
\noindent \VUgras{Tableau 4.} --- L'Age mûr (Le Repas de noce) \dotfill 23

%%K t99-07-1 p
\noindent La Vieillesse (L'Anniversaire) \dotfill 26

%%K t99-08-1 p
\noindent \VUgras{Tableau 5.} --- La Maison et sa construction \dotfill 29

%%K t99-09-1 p
\noindent \VUgras{Tableau 6.} --- La Maison intérieure. --- Les Meubles \dotfill 35

%%K t99-10-1 p
\noindent \VUgras{Tableau 7.} --- Le Village (en hiver) \dotfill 43

%%K t99-11-1 p
\noindent La Maison rustique (au printemps) \dotfill 46

%%K t99-12-1 p
\noindent \VUgras{Tableau 8.} --- La Moisson (en été) \dotfill 49

%%K t99-13-1 p
\noindent La Vendange (en automne) \dotfill 51

%%K t99-14-1 p
\noindent \VUgras{Tableau 9.} --- La Montagne. --- La Ville d'eaux \dotfill 54

%%K t99-15-1 p
\noindent La Forêt. --- La Chasse \dotfill 57

%%K t99-16-1 p
\noindent \VUgras{Tableau 10.} --- La Mer. --- Le Port \dotfill 60

%%K t99-17-1 p
\noindent \VUgras{Tableau 11.} --- La Ville et ses monuments. --- Un\nl
incendie \dotfill 65

%%K t99-18-1 p
\noindent \VUgras{Tableau 12.} --- La Gare. --- Les Chemins de fer \dotfill 69

%%K t99-19-1 p
\noindent \VUgras{Tableau 13.} --- L'Hôtel. --- Le Restaurant. --- Le Café \dotfill 73

%%K t99-20-1 p
\noindent \VUgras{Tableau 14.} --- La Rue. --- Les Commerçants \dotfill 76

%%K t99-21-1 p
\noindent \VUgras{Tableau 15.} --- Le Marché. --- Les Comestibles \dotfill 82

%%K t99-22-1 p
\noindent \VUgras{Tableau 16.} --- Un Grand Magasin. --- Les Jouets \dotfill 86}

\VUsaut{6.0mm}
\VUfilet{20mm}
\end{VUpage}

% --- Feuillet 94 (page blanche, non foliotee) -------------------------
\begin{VUpage}[94]{}
% (page vierge)
\end{VUpage}

% --- Feuillet 95 (acheve d'imprimer, non foliote) ---------------------
\begin{VUpage}[95]{}
%%K t99-apar-1 apar
\VUsaut{60.0mm}
\VUfilet{18mm}
\VUsaut{3.0mm}
\VUcentre{8.0pt}{0}{44872. --- Bordeaux. --- Imp. G. Delmas, rue Saint-Christoly, 10.}
\VUsaut{3.0mm}
\VUfilet{18mm}
\end{VUpage}

% --- Feuillet 96 (4e de couverture, non foliote) -----------------------
%  Publicite « Collection Delmas », en francais, sur mesure de cadre
%  ornemental non reproduite ici (seul le texte l'est).
\begin{VUpage}[96]{}
%%K t99-apar-9 apar
% LA QUATRIEME DE COUVERTURE EST UNE PLANCHE, non une page de texte,
% et un encadrement Art nouveau la borde. On la compose desormais en
% caracteres, comme les feuillets 110 et 113 : \VUkadro dessine le
% cadre (double filet, fleurons d'angle), et chaque ligne prend la
% macro qui convient a son dessin -- \VUtitre/\VUcentre pour les
% lignes d'apparat centrees, \VUlineoG pour les titres « I. --- »,
% « II. --- », « III. --- » et pour les enumerations « 1^o », « 2^o »,
% qui se composent au fer a gauche, a leur chasse naturelle : prises
% a tort dans un alinea justifie, elles s'etiraient d'une marge a
% l'autre, avec des blancs de deux cadratins entre leurs mots.
% Les paragraphes qui remplissent reellement la justification (les
% cinq notices numerotees de la section III, par exemple) restent du
% texte courant et gardent leurs coupures \nl.
\VUkadro{463pt}{%
\VUsaut{5.0mm}
\VUtitre{15.0pt}{60}{Collection Delmas}
\VUsaut{2.0mm}
\VUcentre{9.0pt}{0}{pour l'Enseignement pratique des Langues vivantes}
\VUsaut{2.0mm}
\VUcentre{9.5pt}{40}{PAR L'IMAGE ET LA MÉTHODE DIRECTE}
\VUsaut{1.6mm}
\VUcentre{8.0pt}{0}{\textit{(Conforme aux programmes de Mai 1902)}}
\VUsaut{2.0mm}
\VUornamento{9pt}
\VUsaut{3.0mm}

\VUlineoG{12.0pt}{20}{0mm}{I. --- Tableaux Auxiliaires Delmas}
\VUsaut{1.6mm}

{\fontsize{8.5pt}{10.2pt}\selectfont
\noindent \textsc{Édition portative} (30\texttimes 39 c.), 16 Tableaux, se vend :\par}
\VUsaut{0.8mm}
\VUlineoG{8.5pt}{0}{0mm}{1\textsuperscript{o} En feuilles séparées ;}
\VUlineoG{8.5pt}{0}{0mm}{2\textsuperscript{o} En cahiers : 1\textsuperscript{er} cahier (N\textsuperscript{os} 1 à 6) ;}
\VUlineoG{8.5pt}{0}{0mm}{3\textsuperscript{o} \hphantom{En cahiers} » \hphantom{aaaa} 2\textsuperscript{e} cahier (N\textsuperscript{os} 7 à 16)}

\VUsaut{3.0mm}
\VUtitre{10.5pt}{20}{Éditions spéciales}
\VUsaut{1.2mm}
\VUcentre{8.5pt}{0}{en Anglais, Allemand, Espagnol, Italien, Russe et Français}

\VUsaut{3.0mm}
\VUlineoG{12.0pt}{20}{0mm}{II. --- Tableaux Muraux en Couleurs}
\VUsaut{1.6mm}

{\fontsize{8.5pt}{10.2pt}\selectfont
\noindent (90\texttimes 120 c.), 16 Tableaux, se vend :\par}
\VUsaut{0.8mm}
\VUlineoG{8.5pt}{0}{0mm}{1\textsuperscript{o} En feuilles séparées ;}
\VUlineoG{8.5pt}{0}{0mm}{2\textsuperscript{o} Première série (N\textsuperscript{os} 1 à 6) ;}
\VUlineoG{8.5pt}{0}{0mm}{3\textsuperscript{o} Deuxième et troisième série (N\textsuperscript{os} 7 à 16).}
\VUlineoG{8.5pt}{0}{0mm}{4\textsuperscript{o} Collection complète (N\textsuperscript{os} 1 à 16).}

\VUsaut{3.0mm}
\VUlineoG{12.0pt}{20}{0mm}{III. --- Livrets explicatifs des 16 tableaux}
\VUsaut{1.2mm}
\VUcentre{8.5pt}{0}{(Livres de lecture)}
\VUsaut{2.0mm}

{\fontsize{8.5pt}{10.2pt}\selectfont
\noindent \VUgras{1\textsuperscript{o} Livret explicatif,} français, par M. \textsc{Rochelle,}\nl
professeur au Lycée de Bordeaux.

\VUblancAlinea
\VUgras{2\textsuperscript{o} Sprech-und Leseübungen} nach Delmas' Hilfs\cc
bildern, von E. \textsc{Rochelle,} Lehrer am Gymna\cc
sium v. Bordeaux, und K. \textsc{Körner,} Oberlehrer am\nl
Realgymnasium zu Neunkirchen (Rheinpreussen).

\VUblancAlinea
\VUgras{3\textsuperscript{o} Conversation and Reading-book,} \textit{based on}\nl
\textit{Delmas' auxiliary Pictures,} par Ch. \textsc{Veillet}\cc
\textsc{Aboison,} professeur à Paris.

\VUblancAlinea
\VUgras{4\textsuperscript{o} Libretto spiegativo} (Lettura et Conversazioni),\nl
par M. A. \textsc{Vannier,} professeur à l'École profes\cc
sionnelle de Rouen.

\VUblancAlinea
\VUgras{5\textsuperscript{o} Librillo explicativo} (Lectura y Conversación)\nl
par M. Th. \textsc{Alaux,} prof. au Lycée de Bordeaux.}%
}
\end{VUpage}
